% 脉冲星行星（综述）
% license CCBYSA3
% type Wiki

本文根据 CC-BY-SA 协议转载翻译自维基百科\href{https://en.wikipedia.org/wiki/Pulsar_planet}{相关文章}。

\begin{figure}[ht]
\centering
\includegraphics[width=6cm]{./figures/427823341c42e767.png}
\caption{带有行星的脉冲星艺术概念图} \label{fig_Pupl_1}
\end{figure}
脉冲星行星是指围绕脉冲星运行的行星。1992 年，人们首次在一颗毫秒脉冲星周围发现了这类行星，它们也是人类确认发现的第一批太阳系外行星。脉冲星本身是极其精确的“宇宙时钟”，即使质量很小的行星，也能在脉冲星的观测特性中引起可探测的变化；目前已知质量最小的系外行星，正是一颗脉冲星行星。

这类行星极为罕见，NASA 系外行星档案中仅列出大约半打（六颗左右）。只有一些非常特殊的过程，才能在脉冲星周围形成行星级别的伴星，其中许多被认为是性质十分奇特的天体，例如由伴星部分被破坏后形成的“钻石行星”。此外，脉冲星强烈的辐射以及由电子–正电子对组成的高速粒子风，往往会剥离这些行星的大气层，因此它们几乎不可能成为生命适宜的栖居地。
\subsection{形成}
行星的形成需要原行星盘的存在，而多数理论还要求盘内存在一个“死区”，即没有湍流的区域。在该区域内，微行星体可以形成并不断积累，而不会因向恒星内落而被吞噬$^{[1]}$。与年轻恒星相比，脉冲星具有高得多的光度，因此其辐射会强烈电离盘物质，从而抑制死区的形成$^{[2]}$；电离会触发磁旋转不稳定性，引发湍流并破坏死区$^{[3]}$。因此，若要孕育行星，围绕脉冲星的盘必须具有相当大的质量$^{[4]}$。
目前提出了若干可能形成行星系统的过程$^{[a]}$：
\begin{itemize}
\item “第一代”行星：指在恒星发生超新星爆发并演化为脉冲星之前，就已围绕其运行的行星$^{[6]}$。大质量恒星往往缺乏行星，这可能既源于在极其明亮的恒星周围难以探测行星，也由于强辐射会破坏原行星盘。若行星轨道半径小于约 4 个天文单位，当恒星演化为红巨星或红超巨星时，行星很可能被吞没并毁灭。在超新星爆发过程中，系统会损失大约一半的质量；除非脉冲星在爆发时被“踢出”的方向恰与行星当时的运动方向一致，否则行星往往会与系统解体而脱离。已知的脉冲星行星系统中，没有一个很可能是通过这一过程形成的$^{[7]}$。
\item “第二代”行星：由超新星爆发后回落到脉冲星上的物质形成$^{[6]}$。理论上，这些回落物质的总质量可与典型原行星盘相当$^{[7]}$，但其耗散速度可能过快，不利于行星形成。目前尚无围绕年轻脉冲星的已知行星实例$^{[8][9]}$。
\item “第三代”行星$^{[6]}$：脉冲星与伴星相互作用并将其摧毁，从而形成一个低质量盘。脉冲星可发出高能辐射加热伴星，使其充满并溢出罗希瓣，最终被完全破坏。另一种机制是引力波辐射导致轨道不断收缩，直到伴星（在这些情形中通常为白矮星）解体$^{[8]}$。第三种机制是脉冲星进入一颗更大恒星的包层，使其解体并在脉冲星周围形成一个盘$^{[10][11]}$。由这些过程形成的盘，其质量远大于回落物质形成的盘，因此能够存活更长时间，从而允许行星形成$^{[8]}$。这些盘还富含行星形成所必需的重元素，并且盘物质的一部分会被脉冲星吸积，从而使其自转加快$^{[12]}$。此外，也可能出现另一种情形：一颗较轻的白矮星在与更大质量白矮星的相互作用中被摧毁，轻白矮星形成的碎屑盘孕育行星，而较大的白矮星最终演化为脉冲星$^{[13]}$。
\item 在某些情况下，伴星在与脉冲星的相互作用中被摧毁，但会留下一个行星尺度的残余体$^{[3]}$。这类系统被称为“黑寡妇”系统$^{[14]}$，但并不总被视为脉冲星行星的一种 $^{[15]}$。
\item 最后，还有可能是来自伴星的行星或流浪行星被脉冲星俘获 $^{[16]}$，或者脉冲星与行星原先所围绕的宿主恒星发生并合 $^{[17]}$。后一种过程会形成一个公共包层，该包层最终瓦解并形成一个盘，从中孕育行星 $^{[18]}$。
\end{itemize}
\subsubsection{影响}
不同的形成情景会对行星的组成产生直接影响：由超新星碎屑形成的行星很可能富含金属和放射性同位素$^{[16]}$，并且可能含有大量的水$^{[19]}$；由白矮星解体形成的行星则会富含碳$^{[16]}$，并可能由大量金刚石构成$^{[20]}$；而真正的白矮星碎片本身将具有极高的密度$^{[16]}$。截至 2022 年，在脉冲星周围发现的最常见行星类型是所谓的“钻石行星”，即一种极低质量的白矮星$^{[21]}$。此外，围绕脉冲星的天体还可能包括小行星、彗星以及小行星体$^{[22]}$。更具推测性的设想还包括由奇异物质构成的行星，这类行星可能比普通物质行星更靠近脉冲星运行，并有潜力发射引力波$^{[23]}$。

行星还可以与脉冲星的磁场发生相互作用，产生所谓的“阿尔芬翼”（Alfvén wings）：这是一种围绕行星形成的翼状电流结构，可向行星注入能量$^{[24]}$，并可能产生可被探测到的射电辐射$^{[25]}$。
\subsection{可观测性}
脉冲星是极其精确的“宇宙时钟”$^{[4]}$，其脉冲计时具有高度的规律性。因此，可以通过脉冲星计时变化来探测其周围非常小的天体，最小甚至可达大型小行星的尺度$^{[1]}$。在分析计时信号时，需要对地球和太阳系运动的影响、脉冲星位置估计误差，以及电磁辐射在星际介质中传播时间的变化进行校正。脉冲星会以高度规律的方式自转并逐渐减慢 $^{[4]}$；行星通过对脉冲星施加引力作用，会改变这一模式，从而在脉冲信号中产生多普勒频移$^{[26]}$。理论上，这种技术还可用于探测脉冲星行星的系外卫星$^{[27]}$。然而，该方法也存在局限：脉冲星“突变”（glitch）以及脉冲模式变化有时会模拟出行星存在的信号$^{[28]}$。

最早被发现的系外行星$^{[b]}$（1992 年由 Dale Frail 和 Aleksander Wolszczan 发现）正是围绕 PSR B1257+12 的脉冲星行星$^{[31]}$。这一发现表明，从地球上可以探测到系外行星$^{[32]}$，并由此引发了“系外行星并不罕见”的预期$^{[4]}$。截至 2016 年 $^{[33]}$，已知质量最小的系外行星（PSR B1257+12 A，仅约 $0.02,M_{\oplus}$）也是一颗脉冲星行星$^{[34]}$。

然而，由于体积极小以及脉冲星复杂而特殊的光谱特征，直接成像这类行星实际上极为困难$^{[16]}$。一种潜在的成像方式是观测行星凌日：但对脉冲星行星而言，由于脉冲星本身尺寸极小，行星发生凌日的几率非常低。同时，脉冲星的复杂光谱也使得对行星进行光谱分析变得十分困难。相比之下，行星磁场—脉冲星相互作用以及行星自身的热辐射，更有可能成为获取行星信息的途径 $^{[35]}$。

此外，脉冲星行星还被用于解释某些天文现象，例如软伽马重复暴（soft gamma repeaters）的 X 射线爆发$^{[36]}$。
\subsection{发生率}
截至 2022 年，仅有大约半打$^{[c]}$ 脉冲星行星被确认$^{[11]}$，这意味着其发生率不超过每 200 颗脉冲星中出现 1 个行星系统$^{[d][40]}$。由于对“行星”定义的差异，不同巡天项目给出的统计数量也不尽相同。大规模脉冲星行星搜索的结果同样支持这种极端稀有性$^{[41]}$。大多数行星形成情景要求其前身为质量差异很大的双星系统，且该系统必须在形成脉冲星的超新星爆发中幸存下来；而这两种条件本身都极少同时满足，因此脉冲星行星的形成是一个罕见过程$^{[3]}$。此外，行星及其轨道还必须经受住脉冲星释放的强烈高能辐射（包括 X 射线、伽马射线和高能粒子，即“脉冲星风”）$^{[6]}$。对于通过吸积而“回旋加速”的毫秒脉冲星而言，这一点尤为关键：在其作为 X 射线双星的阶段，所释放的辐射可能会蒸发行星$^{[42]}$。再者，脉冲星的可观测寿命通常只有几百万年，短于行星形成所需的时间，从而进一步限制了我们观测到它们的机会 $^{[43]}$。

基于已知的发生率估计，银河系中可能存在多达一千万颗脉冲星行星$^{[e][46]}$。目前所有已知的脉冲星行星都位于毫秒脉冲星周围$^{[1]}$，这些都是通过从伴星吸积物质而被“加速自转”的年老脉冲星。截至 2015 年，尚未发现围绕年轻脉冲星的行星$^{[47]}$；这是因为年轻脉冲星的脉冲规律性较差，计时误差更大，从而使行星探测更加困难$^{[35]}$。
\subsubsection{已确认的脉冲星行星}
\begin{figure}[ht]
\centering
\includegraphics[width=14.25cm]{./figures/3cd1c1495c4b2471.png}
\caption{} \label{fig_Pupl_2}
\end{figure}
\textbf{M62H}

M62H 是一颗位于蛇夫座的毫秒脉冲星。它位于球状星团 Messier 62 中$^{[62]}$，距离地球约 5 600 秒差距（18 000 光年$^{[63]}$。该脉冲星于 2024 年通过 MeerKAT 射电望远镜被发现$^{[62]}$。M62H 的自转周期为 3.70 毫秒，这意味着它每秒完成约 270 次自转（270 Hz$^{[64]}$。其行星伴星的最小质量为 2.5 个木星质量（$M_J$），在假设脉冲星质量为 1.4 个太阳质量（$M_\odot$）的情况下，其中值质量为 2.83$M_J$。其最小密度为 11 g/cm$^3$。若采用中值质量，则意味着其最大半径约为 48 850 千米（30 350 英里）$^{[65]}$。该行星仅需 0.133 天（3.2 小时）即可完成一次公转，其轨道半径相当于天文单位的 0.49\%，绕 M62H 运行$^{[66]}$。

\textbf{PSR B1257+12}

脉冲星 PSR B1257+12 位于室女座，距离地球约$710^{+43}_{-38}$秒差距$^{[67]}$，其拥有行星的事实于 1992 年通过阿雷西博天文台的观测得到确认$^{[68]}$。该系统包含一颗质量仅为$0.02\pm0.002$ 个地球质量的小型行星，以及两颗质量分别为$4.3\pm0.2$和$3.9\pm0.2$ 个地球质量的超级地球，假设脉冲星质量为 1.4 个太阳质量$^{[69]}$。这些行星极有可能形成于一个原行星盘$^{[1]}$，该盘可能源自伴星的部分破坏$^{[8]}$。计算机模拟表明，该系统至少在十亿年的时间尺度上是稳定的$^{[69]}$，并且系外卫星也有可能在该系统中长期存活$^{[70]}$。该系统在结构上类似于太阳系内区$^{[4]}$；这些行星绕脉冲星运行的距离与水星绕太阳的距离相当，其表面温度也可能处于类似水平$^{[71]}$。关于该系统中存在额外天体的部分报道，可能源于太阳活动造成的扰动$^{[72]}$。

\textbf{PSR J1719−1438}

一颗冥王式行星（cthonian planet）$^{[73]}$ 围绕脉冲星 PSR J1719−1438运行，其质量与木星相当，但半径不足木星的 40\% $^{[i][1]}$。该行星很可能是原伴星在脉冲星强辐射作用下被蒸发后留下的富碳残骸$^{[3]}$，因此常被称为一颗“钻石行星”$^{[j][6]}$。

\textbf{PSR B1620−26}

一颗质量约为$2.5\pm1$ 个木星质量的环双星行星$^{[75]}$，围绕位于球状星团 M4 中的双星系统 PSR B1620−26 运行$^{[4]}$。该双星系统由一颗脉冲星和一颗白矮星组成$^{[1]}$。这颗行星可能是通过引力俘获进入脉冲星轨道的，在球状星团这种高密度环境中，该过程尤为可能发生$^{[16]}$。其年龄可能高达约 126 亿年，使其成为已知最古老的行星$^{[k][76]}$。这一发现也表明，即使在低金属丰度的环境（如球状星团）中，行星同样可以形成$^{[77]}$。

\textbf{PSR J2322−2650}

PSR J2322−2650 似乎拥有一颗质量大致与木星相当的伴星。来自脉冲星的辐射可能将该伴星加热至约 2300 K；在脉冲星附近观测到的一个光源，可能正是这颗行星$^{[78]}$。该脉冲星的光度显著低于许多其他脉冲星，这或许解释了该行星为何能够存活至今$^{[79]}$。利用 JWST NIRSpec 的观测发现，这颗行星的大气层富含分子态碳（如 C$_3$、C$_2$），并存在强烈的向西风$^{[80]}$。
\subsubsection{碎屑盘与前身天体}
对脉冲星 PSR B1937+21 和 PSR J0738−4042 的计时变化，可能反映了这些脉冲星周围存在一条小行星带$^{[l]}$。此外，有学者提出，小行星或彗星与脉冲星发生碰撞，可能是快速射电暴$^{[m]}$、伽马射线暴 GRB 101225A$^{[6]}$ 以及其他类型脉冲星变异现象的成因之一$^{[84]}$。目前尚未在脉冲星周围发现已知的碎屑盘，不过有研究认为磁星 4U 0142+61 和 1E 2259+586$^{[n]}$ 可能拥有此类结构 $^{[2]}$。

白矮星—脉冲星双星系统 PSR J0348+0432 可能是一个未来有潜力形成脉冲星行星的系统$^{[86]}$。此外，也有人提出，在脉冲星 Geminga 周围存在的尘埃云，可能是行星形成的前身结构$^{[87]}$。
\subsubsection{候选体}
早期曾有一些关于脉冲星行星的报道，后来被撤回或被认为证据不足$^{[88]}$。例如，1991 年所谓在 PSR B1829−10 周围发现行星的“发现”，后来被证明只是由地球运动引起的观测伪影$^{[4]}$。自 1979 年以来，关于脉冲星 PSR B0329+54 是否存在行星的问题一直存在争议，截至 2017 年仍未得到解决$^{[89]}$。而 PSR B1828−11 已被明确证实，其表现出的现象源于磁层活动，会模拟出类似行星的信号，但实际上并不存在行星$^{[90]}$。同样，曾提出的围绕脉冲星 Geminga 的行星候选体，后来也被归因于计时噪声$^{[87]}$。
\begin{figure}[ht]
\centering
\includegraphics[width=14.25cm]{./figures/9dea5ff9cc1ac3a9.png}
\caption{} \label{fig_Pupl_3}
\end{figure}
\subsection{宜居性}
脉冲星发射的辐射谱与普通恒星截然不同：几乎不产生可见光或红外辐射，但会释放大量电离辐射$^{[46]}$以及电子–正电子对，这些粒子由脉冲星在自转过程中其强磁场所产生。此外，脉冲星诞生前遗留下来的热量、其自身辐射对两极的加热以及物质吸积过程，都会驱动热辐射和中微子的发射$^{[100]}$。电子–正电子对和 X 射线会被行星大气吸收并加热，从而引发强烈的大气逃逸，甚至可能将大气彻底剥离$^{[101]}$。若行星具备自身的磁场，可能在一定程度上缓解电子–正电子对的影响$^{[102]}$。

传统上，宜居性通过行星的平衡温度来定义，而平衡温度取决于其接收的辐射通量；若行星表面能够存在液态水，则被称为“宜居”$^{[103]}$。不过，即便外部能量输入很少，一些行星也可能在地下孕育生命$^{[104]}$。由于脉冲星体积很小，其辐射总量并不大，因此其宜居带往往会非常靠近恒星本身，以至于潮汐效应可能直接摧毁行星$^{[105]}$。此外，对于某一具体脉冲星，其实际辐射强度以及有多少辐射能够到达假想行星表面，往往难以准确判定；在已知的脉冲星行星中，只有 PSR B1257+12 系统中的行星接近宜居带$^{[106]}$，而截至 2015 年，没有任何已知的脉冲星行星被认为可能是宜居的$^{[4][39]}$。额外的热源还可能包括在导致脉冲星形成的超新星过程中生成的放射性同位素（如 钾-40）$^{[19]}$，以及近轨道行星所经历的潮汐加热$^{[107]}$。来自伴星等外部天体的辐射，也会增加行星的能量收支$^{[73]}$。
\subsection{参见}
\begin{itemize}
\item 行星列表
\item 具有原行星盘的恒星列表
\item 安德鲁·莱恩（Andrew Lyne）
\end{itemize}